\scanpage{36}
[ $\mathcal X$ ens ordonné, $A\subset\mathcal X$, $B$ l'ens des $b\in\mathcal X$ qui majorent $A$

\deleted{\unsure{$\exists$}} Par définition, $\mathrm{Sup}_{\mathcal X}A$ existe \unsure{sss|ssi} $B$ a un plus petit élément $\beta$,
et on pose \deleted{$\unsure{\beta=}$} $\mathrm{Sup}_{\mathcal X}A=\beta$. Si p. ex. $A$ est vide, $B=\mathcal X$, et
$\mathrm{Sup}_{\mathcal X}\varnothing$ existe sss $\mathcal X$ a un plus petit élément (qui
est alors $\mathrm{Sup}_{\mathcal X}\varnothing$). On \unsure{définit|definit} dualement $\mathrm{Inf}_{\mathcal X}A$\ldots

\uline{Prop. 1} Si $A$ a un plus grand élément $\alpha$, il a un sup. et c'est $\alpha$. (Trivial)

\uline{Prop} 2 Pour que $\mathrm{Sup}_{\mathcal X}A$ existe, i.e $B$ ait un plus petit
él., il f et s que $\mathrm{Inf}_{\mathcal X}B$ existe, et alors
$\mathrm{Inf}_{\mathcal X}B=\mathrm{Sup}_{\mathcal X}A$. c'est évidemment néc \deleted{\unsure{q}} d'après prop. 1. Suff?

Soit $\widetilde A$ l'ens des $a\in\mathcal X$ qui minorent $B$. Donc
$A\subset\widetilde A\unsure{\prec|<}B$. Par hypothèse, ($B$ admet un Inf) $\widetilde A$ a un
plus grand él. $\alpha$, lequel majore donc $\widetilde A$ et a fortiori
$A$, donc $\alpha\in B$, et étant dans $\unsure{\bar A|\widetilde A}$ il minore $B$, donc il
est un plus petit él. de $B$, cqfd.

\uline{\deleted{\unsure{Prop}} Cor} Si tte partie de $\mathcal X$ admet un Inf, alors tte
partie de $\mathcal X$ admet un Sup, et réciproquement.

Si tte partie \deleted{non vide} minorée \inserted{(resp. non vide)} de $\mathcal X$ admet un
Inf, alors tte partie non vide (resp. majorée) de
$\mathcal X$ admet un Sup, et réciproquement. Donc
pour que tte partie min. non vide de $\mathcal X$ admette
un Inf, il f et s que tte partie majorée non
vide de $\mathcal X$ admette un Sup. \hfill]
% archive sheet number: 35
